% GENERATED FILE. Edit canonical lessons, then run npm run generate:books.
% canonical-source-hash: fnv1a64:53e555b1ba5d269f
% canonical-lessons: FR-C524-version-originale, FR-C524-projeter, FR-C524-doublage, FR-C524-realisateur, FR-C524-tournage

\chapter{En version originale}
\label{ch:fr-en-version-originale}
\hlchaptermodality{524}

\begin{chapteropening}
\textbf{By the end of this chapter:} \emph{I can talk about how a film is made and shown.}

Complete the last lesson of chapter 524: I can talk about how a film is made and shown.
\end{chapteropening}

\section[la version originale]{la version originale — original-language version}

\label{lesson:FR-C524-version-originale}



\begin{quote}
Before the new one: say the French for (streaming) platform, then the French for thrilling.
\end{quote}

\subsection*{You'll want to know: la version originale}
\textbf{la version originale} — \textquotedblleft{}original-language version\textquotedblright{}.

\begin{quote}
Je préfère voir les films en version originale.
\end{quote}

\begin{quote}
\emph{I prefer to watch films in the original language.}
\end{quote}

Often \emph{la VO}; \emph{la VOST} has subtitles; \emph{la VF} is the French dub.

\begin{cousinweb}
\emph{la version originale} joins \emph{version}, from Medieval Latin \emph{versiō} \textquotedblleft{}a turning\textquotedblright{}, from \emph{vertere} \textquotedblleft{}to turn\textquotedblright{}, and \emph{originale}, from Latin \emph{orīgō} \textquotedblleft{}beginning\textquotedblright{}, from \emph{orīrī} \textquotedblleft{}to rise\textquotedblright{} (Lewis \& Short).
\end{cousinweb}

\begin{grammarlens}[title={What you've built}]
One more word for telling, explaining and describing at B1.
\end{grammarlens}

\subsection*{Your turn}
\begin{itemize}
\raggedright
  \item \emph{Say it:} \emph{la version originale}
  \item \emph{Say it:} \emph{la version originale}, once more
  \item \emph{Say it:} a sentence of your own with \emph{la version originale}
  \item \emph{Recall:} say the French for (streaming) platform, then the French for thrilling, then say \emph{la version originale} again
\end{itemize}

\begin{tcolorbox}[breakable,colback=teal!4,colframe=teal!35!black,title={Before you move on}]
What does la version originale mean? (\textquotedblleft{}Original-language version\textquotedblright{}.) Say it once more.
\end{tcolorbox}
\section[projeter]{projeter — to screen (a film); to plan; to project}

\label{lesson:FR-C524-projeter}



\begin{quote}
Before the new one: say the French for thrilling, then the French for original-language version.
\end{quote}

\subsection*{You'll want to know: projeter}
\textbf{projeter} — \textquotedblleft{}to screen (a film); to plan; to project\textquotedblright{}.

\begin{quote}
Le cinéma du quartier projette des films en version originale.
\end{quote}

\begin{quote}
\emph{The local cinema screens films in the original language.}
\end{quote}

Spelling: \emph{je projette}, \emph{nous projetons}. Also to plan: \emph{projeter un voyage}.

\begin{cousinweb}
\emph{projeter} comes from Old French \emph{porjeter} \textquotedblleft{}to throw forward\textquotedblright{}, built on \emph{jeter} \textquotedblleft{}to throw\textquotedblright{}, which goes back to Latin \emph{iactāre}, from \emph{iacere} \textquotedblleft{}to throw\textquotedblright{} (Lewis \& Short).
\end{cousinweb}

\begin{grammarlens}[title={What you've built}]
One more word for telling, explaining and describing at B1.
\end{grammarlens}

\subsection*{Your turn}
\begin{itemize}
\raggedright
  \item \emph{Say it:} \emph{projeter}
  \item \emph{Say it:} \emph{projeter}, once more
  \item \emph{Say it:} a sentence of your own with \emph{projeter}
  \item \emph{Recall:} say the French for thrilling, then the French for original-language version, then say \emph{projeter} again
\end{itemize}

\begin{tcolorbox}[breakable,colback=teal!4,colframe=teal!35!black,title={Before you move on}]
What does projeter mean? (\textquotedblleft{}To screen (a film); to plan; to project\textquotedblright{}.) Say it once more.
\end{tcolorbox}
\section[le doublage]{le doublage — dubbing}

\label{lesson:FR-C524-doublage}



\begin{quote}
Before the new one: say the French for original-language version, then the French for to screen.
\end{quote}

\subsection*{You'll want to know: le doublage}
\textbf{le doublage} — \textquotedblleft{}dubbing\textquotedblright{}.

\begin{quote}
Le doublage de ce dessin animé est très réussi.
\end{quote}

\begin{quote}
\emph{The dubbing of this cartoon is very good.}
\end{quote}

Built on \emph{doubler} \textquotedblleft{}to dub\textquotedblright{}.

\begin{cousinweb}
\emph{le doublage} is built on \emph{doubler} with \emph{-age}; \emph{doubler} comes from \emph{double}, from Latin \emph{duplus} \textquotedblleft{}twofold\textquotedblright{}, built on \emph{duo} \textquotedblleft{}two\textquotedblright{} (Lewis \& Short).
\end{cousinweb}

\begin{grammarlens}[title={What you've built}]
One more word for telling, explaining and describing at B1.
\end{grammarlens}

\subsection*{Your turn}
\begin{itemize}
\raggedright
  \item \emph{Say it:} \emph{le doublage}
  \item \emph{Say it:} \emph{le doublage}, once more
  \item \emph{Say it:} a sentence of your own with \emph{le doublage}
  \item \emph{Recall:} say the French for original-language version, then the French for to screen, then say \emph{le doublage} again
\end{itemize}

\begin{tcolorbox}[breakable,colback=teal!4,colframe=teal!35!black,title={Before you move on}]
What does le doublage mean? (\textquotedblleft{}Dubbing\textquotedblright{}.) Say it once more.
\end{tcolorbox}
\section[le réalisateur]{le réalisateur — (film) director}

\label{lesson:FR-C524-realisateur}



\begin{quote}
Before the new one: say the French for to screen, then the French for dubbing.
\end{quote}

\subsection*{You'll want to know: le réalisateur}
\textbf{le réalisateur} — \textquotedblleft{}(film) director\textquotedblright{}.

\begin{quote}
Le réalisateur a présenté son film au public.
\end{quote}

\begin{quote}
\emph{The director presented his film to the audience.}
\end{quote}

A woman is \emph{la réalisatrice}. In a theatre, \emph{le metteur en scène}.

\begin{cousinweb}
\emph{le réalisateur} is built on \emph{réaliser} \textquotedblleft{}to make real\textquotedblright{}, from \emph{réel}, from Late Latin \emph{reālis}, built on \emph{rēs} \textquotedblleft{}a thing\textquotedblright{} (Lewis \& Short).
\end{cousinweb}

\begin{grammarlens}[title={What you've built}]
One more word for telling, explaining and describing at B1.
\end{grammarlens}

\subsection*{Your turn}
\begin{itemize}
\raggedright
  \item \emph{Say it:} \emph{le réalisateur}
  \item \emph{Say it:} \emph{le réalisateur}, once more
  \item \emph{Say it:} a sentence of your own with \emph{le réalisateur}
  \item \emph{Recall:} say the French for to screen, then the French for dubbing, then say \emph{le réalisateur} again
\end{itemize}

\begin{tcolorbox}[breakable,colback=teal!4,colframe=teal!35!black,title={Before you move on}]
What does le réalisateur mean? (\textquotedblleft{}(film) director\textquotedblright{}.) Say it once more.
\end{tcolorbox}
\section[le tournage]{le tournage — shooting, filming (of a film)}

\label{lesson:FR-C524-tournage}



\begin{quote}
Before the new one: say the French for dubbing, then the French for (film) director.
\end{quote}

\subsection*{You'll want to know: le tournage}
\textbf{le tournage} — \textquotedblleft{}shooting, filming (of a film)\textquotedblright{}.

\begin{quote}
Le tournage a duré trois mois.
\end{quote}

\begin{quote}
\emph{The filming lasted three months.}
\end{quote}

\emph{Tourner un film} — to shoot a film.

\begin{cousinweb}
\emph{le tournage} is built on \emph{tourner} with \emph{-age} (EtymDB); \emph{tourner} comes from Latin \emph{tornāre} \textquotedblleft{}to turn on a lathe\textquotedblright{} (Lewis \& Short), and early cameras were cranked by hand.
\end{cousinweb}

\begin{grammarlens}[title={What you've built}]
One more word for telling, explaining and describing at B1.
\end{grammarlens}

\subsection*{Your turn}
\begin{itemize}
\raggedright
  \item \emph{Say it:} \emph{le tournage}
  \item \emph{Say it:} \emph{le tournage}, once more
  \item \emph{Say it:} a sentence of your own with \emph{le tournage}
  \item \emph{Recall:} say the French for dubbing, then the French for (film) director, then say \emph{le tournage} again
\end{itemize}

\begin{tcolorbox}[breakable,colback=teal!4,colframe=teal!35!black,title={Before you move on}]
What does le tournage mean? (\textquotedblleft{}Shooting, filming (of a film)\textquotedblright{}.) Say it once more.
\end{tcolorbox}
